\documentclass[11pt]{article}
\usepackage{fontspec}
\setmainfont{Times New Roman}
\setsansfont{Arial}
\setmonofont{Consolas}
\usepackage{amsmath,amssymb}
\usepackage{geometry}
\usepackage{microtype}
\geometry{a4paper,margin=3cm}
\setlength{\parindent}{1.25em}
\setlength{\parskip}{0pt}
\emergencystretch=10em
\allowdisplaybreaks
\providecommand{\Kfield}{\mathfrak{K}}
\providecommand{\Ssys}{\mathfrak{S}}
\providecommand{\Jdom}{\mathfrak{J}}
\providecommand{\Gdom}{\mathfrak{G}}
\providecommand{\Lsys}{\mathfrak{L}}
\providecommand{\Omegaint}{[\Omega]}
\providecommand{\eps}{\varepsilon}

\begin{document}

\section*{\S\ 9. Відносно цілі області $\Gdom_{n\rho}$ з поліномів.}

Далі ми розглядаємо спеціальні інтегральні області з поліномів, які утворюють проміжний рівень між інтегральною областю та полем.

Інтегральна область з поліномів $\Gdom_{n\rho}$ називається ``відносно цілою областю'', якщо виконується така умова:

\emph{4a) $\Gdom_{n\rho}$ разом із $F(x)$ та $G(x)$ --- де $G(x)\ne0$ --- містить частку $F(x):G(x)$ тоді й лише тоді, коли ця частка є цілою за невизначеними, тобто є поліномом.}

Спершу треба довести тотожність відносно цілих областей із сукупністю поліномів деякого поля $\Kfield_{n\sigma}$ ($\sigma\ge \rho$) раціональних функцій.

Нехай $\Kfield_{n\rho}$ є найменшим полем, що містить $\Gdom_{n\rho}$. Для кожної функції $f(x)$ з $\Kfield_{n\rho}$ існує подання як частки двох поліномів із $\Gdom_{n\rho}$:
\[
  f(x)=F_1(x):F_2(x),
\]
де також допускаються поліноми нульового степеня. Щойно $f(x)$ стає поліномом, вона за умовою 4a) належить до $\Gdom_{n\rho}$; а оскільки $\Gdom_{n\rho}$ не може містити нічого понад поліноми, то \emph{кожна відносно ціла область $\Gdom_{n\rho}$ є тотожною сукупності поліномів найменшого поля, що її містить.}

Навпаки, нехай $\Kfield_{n\sigma}$ є довільним полем раціональних функцій (отже, не обов'язково найменшим полем, що містить задану систему), і нехай $\Jdom$ є сукупністю поліномів, які містяться в $\Kfield_{n\sigma}$. $\Jdom$ задовольняє умови 1, 2, 3 інтегральної області та умову 4a), отже є відносно цілою областю; тобто \emph{сукупність поліномів, що містяться в полі $\Kfield_{n\sigma}$, утворює відносно цілу область $\Gdom_{n\rho}$ ($\rho\le\sigma$).}\footnote{Можливий також випадок $\rho=0$; тобто сукупність поліномів із $\Kfield_{n\sigma}$ складається з елементів області коефіцієнтів $\Omega$.}

Крім того, ми покажемо тотожність відносно цілих областей з уведеними Гільбертом ``інтегральними областями з відносно цілих функцій''.\footnote{D. Hilbert: Mathematische Probleme. Доповідь, Париж 1900. Problem 14 (Göttinger Nachrichten 1900).}

Нехай
\begin{equation}
  G_1(x),\ldots,G_k(x)
\tag{1}
\end{equation}
є раціональною базою $\Gdom_{n\rho}$. За умовами 1, 2, 3, 4a кожна раціональна комбінація поліномів (1), яка є цілою за невизначеними, тобто стає поліномом, належить до $\Gdom_{n\rho}$. Оскільки, навпаки, кожний поліном із $\Gdom_{n\rho}$ за поняттям раціональної бази можна раціонально виразити через поліноми (1), то $\Gdom_{n\rho}$ є \emph{тотожною сукупності тих раціональних комбінацій поліномів (1), які є цілими за невизначеними, тобто є поліномами.} Отже, ми маємо означення Гільберта, що виходить зі спеціальної раціональної бази, тоді як наше означення використовує лише абстрактні властивості області.

Для відносно цілих областей означення спільного дільника, дане в \S\ 8 для загальних інтегральних областей, спрощується:

Нехай $F(x),G(x)$ --- два поліноми з $\Gdom_{n\rho}$, а $L(x)$ --- спільний дільник цих поліномів. Тоді $L(x)$ називається \emph{спільним дільником у $\Gdom_{n\rho}$} тоді й лише тоді, коли $L(x)$ належить до $\Gdom_{n\rho}$.

Справді, тоді належність часток $F(x):L(x)$ та $G(x):L(x)$ випливає з умови 4a).

Ще одна істотна властивість відносно цілих областей полягає в тому, що поняття ``алгебраїчно цілого відносно $\Gdom_{n\rho}$'' можна означати за допомогою довільного рівняння точно так само, як це відомо для алгебраїчних числових полів або для поля $\Omega(x_1,\ldots,x_n)$.

\emph{Означення V: Величина $\xi$ називається алгебраїчно цілою відносно $\Gdom_{n\rho}$ --- або також відносно алгебраїчно цілою ---, якщо вона задовольняє деяке рівняння, старший коефіцієнт якого є одиницею, а решта коефіцієнтів є поліномами з $\Gdom_{n\rho}$.}

Справді, нехай відносно алгебраїчно ціла величина $\xi$ задовольняє, наприклад, рівняння
\[
  L(z)=z^\lambda+G_1(x)z^{\lambda-1}+\cdots+G_\lambda(x)=0.
\]
$L(z)$ ділиться на цілу функцію
\[
  M(z)=z^\mu+H_1(x)z^{\mu-1}+\cdots+H_\mu(x),
\]
незвідну відносно найменшого поля $\Kfield_{n\rho}$, що містить $\Gdom_{n\rho}$. Але з цієї подільності, за відомим узагальненням теореми Гауса,\footnote{Пор. наприклад J. König: Einleitung in die allgemeine Theorie der algebraischen Größen (Leipzig, B.~G. Teubner, 1903), Kap. IX, \S\ 2, або Weber (мале видання), \S\ 20.} випливає, що раціональні функції $H_1(x),\ldots,H_\mu(x)$ є поліномами з $\Kfield_{n\rho}$, а отже належать до $\Gdom_{n\rho}$. Так означення відносно алгебраїчно цілої величини $\xi$ отримується з незвідного рівняння. Зокрема, кожний поліном $F(x)$ із $\Gdom_{n\rho}$ також є алгебраїчно цілим відносно $\Gdom_{n\rho}$, бо він задовольняє рівняння $z-F(x)=0$.

Далі зазначимо, що для інтегральної області $\Jdom_{n\rho}$ з поліномів, аналогічно до найменшого поля, що її містить, визначена також найменша відносно ціла область $\Gdom_{n\rho}$, яка її містить, такою вимогою:

$\Gdom_{n\rho}$ разом із $\Jdom_{n\rho}$ містить усі поліноми $H(x)$, і лише їх, які можна подати як частки двох поліномів із $\Jdom_{n\rho}$.

Із цього означення далі випливає: з кожної інволюційної форми та інволюційної бази $\Jdom_{n\rho}$ відповідна форма та база для $\Gdom_{n\rho}$ виникають так, що треба щонайбільше відокремити спільний дільник у $\Gdom_{n\rho}$ усіх поліномів бази.

Адже $\Jdom_{n\rho}$ та $\Gdom_{n\rho}$ мають те саме найменше поле $\Kfield_{n\rho}$, що їх містить; і кожна інволюційна форма $\Jdom_{n\rho}$ виникає з інволюційної форми $\Kfield_{n\rho}$ множенням на поліном із $\Jdom_{n\rho}$, а тому має коефіцієнти з $\Gdom_{n\rho}$, які, однак, можуть мати спільний дільник у $\Gdom_{n\rho}$.\footnote{Наприклад, для інтегральної області $\Jdom_{22}$, що складається з усіх цілих раціональних комбінацій $x_1^2$ та $x_1x_2$, як інволюційна форма отримується:
\[
\Psi(z,u)=x_1^2z^2-x_1^4u_1^2-2x_1^3x_2u_1u_2-x_1^2x_2^2u_2^2.
\]
Оскільки $x_1^2$ належить до $\Jdom_{22}$, а тим більше до найменшої відносно цілої області $\Gdom_{22}$, що її містить, і, отже, є спільним дільником у $\Gdom_{22}$, то як інволюційна форма $\Gdom_{22}$ отримується:
\[
\Psi(z,u)=z^2-x_1^2u_1^2-2x_1x_2u_1u_2-x_2^2u_2^2.
\]}

Нарешті зазначимо, що \emph{область відображення} відносно цілої області $\Gdom_{n\rho}$, яка за \S\ 7 знову є інтегральною областю $\Jdom^*_{\rho\rho}$ з поліномів, не обов'язково сама мусить бути відносно цілою областю. Бо якщо $F(x)$ та $G(x)$ --- два поліноми з $\Gdom_{n\rho}$, частка яких не є поліномом і тому не належить до $\Gdom_{n\rho}$, то частка $F^*(x):G^*(x)$ також не належить до $\Jdom^*_{\rho\rho}$. Однак ця частка, що виникає спеціалізацією деяких невизначених, цілком може бути поліномом; тоді для $\Jdom^*_{\rho\rho}$ умова 4a) не виконується.\footnote{Нехай, наприклад, $\Gdom_{22}$ складається з усіх поліномів, які є раціональними комбінаціями $F=x_1^2x_2$ та $G=x_1^2(1+x_3)$. Допустима спеціалізація $x_3=0$ дає $F^*:G^*=x_2$, тоді як $F:G$ не є поліномом. Отже, $\Jdom^*_{22}$ не є відносно цілою областю.}

\end{document}
